%%%%%%%%%%%%%%%%%%%%%%%%%%%%%%%%%%%%%%%%%%%%%%%%%%%%%%%%%%%%%%%%%%%%%%%
%% Template for Presentations of Master- and Semester projects 
%% at ISI
%% 
%% 24.10.2006, Mich�le Wigger
%%
%%%%%%%%%%%%%%%%%%%%%%%%%%%%%%%%%%%%%%%%%%%%%%%%%%%%%%%%%%%%%%%%%%%%%

\documentclass{beamer}

% to create a PDF:
% ---------------
%
% latex presentation.tex
% latex presentation.tex
% dvips -P pdf -o presentation.ps presentation.dvi
% ps2pdf presentation.ps
%
%
% to create nice two-up handouts:
% ------------------------------
% 
% as header use
%
% \documentclass[dvips,handout]{beamer}
%
% and make sure all pictures are available as in EPS format.
%
% Then:
%
% latex presentation.tex
% dvips -P pdf -ta4 -o presentation.tmp.ps presentation.dvi
% psnup -4 -b1cm -m1cm -W128mm -H96mm -pa4 presentation.tmp.ps presentation.handout.ps
%
%

\mode<presentation>
{
%\usetheme{Singapore}
   \useoutertheme[subsection=false]{smoothbars}
   \useoutertheme[subsection=false,footline=authortitle]{miniframes}
   \useinnertheme[shadow=false]{rounded}
   \usecolortheme{orchid}
   \usecolortheme{whale}
}

\setbeamersize{text margin left=1em,text margin right=1em}
\usefonttheme[onlymath]{serif}
\setbeamertemplate{navigation symbols}{}


% Theorems
%
% redefines the beginning of a theorem to omit parentheses when
% an optional argument is given.
%
% Thus
%   \begin{theorem}[1] ... \end{theorem}
% results in
%   Theorem 1
% ( instead of Theorem (1) )
\setbeamertemplate{theorem begin}{%
  \begin{\inserttheoremblockenv}{%
      \inserttheoremname
      \ifx\inserttheoremaddition\@empty\else\ \inserttheoremaddition \fi%
    }%
}

% redefine footer
\setbeamertemplate{footline}{%
  \begin{beamercolorbox}[colsep=1.5pt]{upper separation line foot}
  \end{beamercolorbox}
 \begin{beamercolorbox}[ht=2.5ex,dp=1.125ex,
 leftskip=.3cm,rightskip=.3cm plus1fil]
{navigation symbols}%title in head/foot%}
 %%%
%%%   \leavevmode{%
 %%%     \insertslidenavigationsymbol%
 %%%     \insertframenavigationsymbol%
%%%    }%
    \hfill
   {\insertframenumber / \inserttotalframenumber}
  \end{beamercolorbox}%
  \begin{beamercolorbox}[colsep=1.5pt]{lower separation line foot}
  \end{beamercolorbox}
}

% redefine header
\setbeamertemplate{headline}{%
  \begin{beamercolorbox}[colsep=1.5pt]{upper separation line head}
  \end{beamercolorbox}
  \begin{beamercolorbox}{section in head/foot}
    \vskip4pt\insertsectionnavigationhorizontal{\paperwidth}{}{}\vskip4pt
  \end{beamercolorbox}%
  \begin{beamercolorbox}[colsep=1.5pt]{lower separation line head}
  \end{beamercolorbox}
}


\usepackage[english]{babel}
\usepackage[latin1]{inputenc}
\usepackage{amsmath,amssymb,mathrsfs}
\usepackage{epstopdf}				% Converts eps-files into pdf files on the fly

\title[Phase-Locked Loops] % (optional, use only with long paper titles)
{Phase-Locked Loops}

\author{Daniel Gilgen and Fabio Marti}
\institute{ETH Zurich\\
\vspace{0.5 cm}
dagilgen@ethz.ch \hspace{1 cm} martif@ethz.ch % \\[4cm]
\vspace{1.4 cm}


}
 \titlegraphic{
 \includegraphics[height=10mm]{ethlogo.eps} \hfill 
 \includegraphics[height=8mm]{isilogo_left.eps}}

\date{\today}


\newcommand{\subsectitle}{}
\newcommand{\subsectiontitle}[1]{\renewcommand{\subsectitle}{#1}}
\newcommand{\vsp}[1]{\vspace{\stretch{#1}}}









%%%%%%%%%%%%%%%%%%%%%%%%%%%%%%%%%%%%%%%%%%%%%%%%%%%%%%%%%%%%%%%%%%%%%%%%%%%%%
%%% Create a title-slide
%%%%%%%%%%%%%%%%%%%%%%%%%%%%%%%%%%%%%%%%%%%%%%%%%%%%%%%%%%%%%%%%%%%%%%%%%%%%%

\begin{document}

\begin{frame}[plain]
  \titlepage
\centering
\end{frame}

%%%%%%%%%%%%%%%%%%%%%%%%%%%%%%%%%%%%%%%%%%%%%%%%%%%%%%%%%%%%%%%%%%%%%%%%%%%%%
%%% Slide-0: Outline of the Presentation
%%%%%%%%%%%%%%%%%%%%%%%%%%%%%%%%%%%%%%%%%%%%%%%%%%%%%%%%%%%%%%%%%%%%%%%%%%%%
\begin{frame}
\frametitle{Outline}
\tableofcontents
\end{frame}

%%%%%%%%%%%%%%%%%%%%%%%%%%%%%%%%%%%%%%%%%%%%%%%%%%%%%%%%%%%%%%%%%%%%%%%%%%%%%
%%% Slide-1
%%%%%%%%%%%%%%%%%%%%%%%%%%%%%%%%%%%%%%%%%%%%%%%%%%%%%%%%%%%%%%%%%%%%%%%%%%%%
\section{Listing Items} %what appears in the header


\subsectiontitle{Standard List} % title on the slide

\subsection{\subsectitle}

\begin{frame}
  
  \frametitle{\subsectitle}

\begin{itemize}
\item Use the itemization ...
\item ... as in the normal latex file
\end{itemize}

\end{frame}

%%%%%%%%%%%%%%%%%%%%%%%%%%%%%%%%%%%%%%%%%%%%%%%%%%%%%%%%%%%%%%%%%%%%%%%%%%%%%
%%% Slide-2
%%%%%%%%%%%%%%%%%%%%%%%%%%%%%%%%%%%%%%%%%%%%%%%%%%%%%%%%%%%%%%%%%%%%%%%%%%%%

\subsectiontitle{Uncovering Items} % title on the slide

\subsection{\subsectitle}

\begin{frame}
  
  \frametitle{\subsectitle}

\begin{itemize}
\item  \uncover<1>{$\backslash$uncover$<1>$ shows the item and the text
     only in the first overlay } 
\item \only<2>{$\backslash$only$<2>$ shows the item always, but the text
    only in the second overlay}  \only<3-4>{ $\backslash$only$<3-4>$
    replaces the text from $\backslash$only$<2>$ in overlays 3 and 4.}
\item \uncover<5-7>{$\backslash$uncover$<5-7>$ keeps the space empty but
  shows the text and the item only in overlays 5 to 7}
\item  This text is showed in all overlays, but starting from overlay
  8, parts of it are {\color<8->[rgb]{1,0,0} red!} 
\item \uncover<9>{ The red color in the above item in the overlays starting from 8 is
  achieved with the command: $\backslash$color$<8->[$rgb$]\{1,0,0\}$}
\end{itemize}

\end{frame}

%%%%%%%%%%%%%%%%%%%%%%%%%%%%%%%%%%%%%%%%%%%%%%%%%%%%%%%%%%%%%%%%%%%%%%%%%%%
%%% Slide-3
%%%%%%%%%%%%%%%%%%%%%%%%%%%%%%%%%%%%%%%%%%%%%%%%%%%%%%%%%%%%%%%%%%%%%%%%%%%%
\section{Including Figures}
 
   \subsectiontitle{Including a Simple Figure}

   \subsection{\subsectitle}
       \begin{frame}
         \frametitle{\subsectitle}

           \begin{center}
             \includegraphics[width=0.5\textwidth]{isilogo_left.eps}
           \end{center}

         You can include eps-files. If you are using pdflatex, you
       have to include png, jpg, or pdf files.
       
\end{frame}
%%%%%%%%%%%%%%%%%%%%%%%%%%%%%%%%%%%%%%%%%%%%%%%%%%%%%%%%%%%%%%%%%%%%%%%%%%%
%%% Slide-4
%%%%%%%%%%%%%%%%%%%%%%%%%%%%%%%%%%%%%%%%%%%%%%%%%%%%%%%%%%%%%%%%%%%%%%%%%%%%
 
   \subsectiontitle{Including Figure with Columns Environment}
   \subsection{\subsectitle}
       \begin{frame}
         \frametitle{\subsectitle}

           \begin{columns}[c] 
             \column{.4\textwidth}   %specifies the size of the left column
                 \begin{itemize}
                    \item The slide is split into two columns
                    \item The left column contains text, but can
                    contain anything, also boxes, another picture, ....
                    \item The right column here contains a picture
                 \end{itemize}        


             \column{.55\textwidth} %specifies the size of the right column
               \includegraphics[width=\columnwidth]{isilogo_left.eps}
           \end{columns}
      \end{frame}     
       

%%%%%%%%%%%%%%%%%%%%%%%%%%%%%%%%%%%%%%%%%%%%%%%%%%%%%%%%%%%%%%%%%%%%%%%%%
%%% Slide-5
%%%%%%%%%%%%%%%%%%%%%%%%%%%%%%%%%%%%%%%%%%%%%%%%%%%%%%%%%%%%%%%%%%%%%%%%%
\section{Emphasizing, Theorems, ...}



 \subsectiontitle{Emphasizing Important Subjects with Blocks}
 
       \subsection{\subsectitle}
       \begin{frame}
         \frametitle{\subsectitle}
        \begin{block}{Title of the block}
          With the block environment you can emphasize results, observations,
          statements, ...
          \begin{itemize}
            \item All standard environments as itemizations are possible
          \end{itemize}
        \end{block}
      
\end{frame}
%%%%%%%%%%%%%%%%%%%%%%%%%%%%%%%%%%%%%%%%%%%%%%%%%%%%%%%%%%%%%%%%%%%%%%%%%
%%% Slide-6
%%%%%%%%%%%%%%%%%%%%%%%%%%%%%%%%%%%%%%%%%%%%%%%%%%%%%%%%%%%%%%%%%%%%%%%%%


 \subsectiontitle{Stating Theorems, Lemmas, Definitions,... }
       \subsection{\subsectitle}
       \begin{frame}
         \frametitle{\subsectitle}

        \begin{definition}[3]          Here you can state the definition. As in the block
          environment all standard environments can be included.
        \end{definition}

        \begin{theorem}[1]
          Instead of theorem you can use lemma, corollary, ... 
        \end{theorem}
      \end{frame}


%%%%%%%%%%%%%%%%%%%%%%%%%%%%%%%%%%%%%%%%%%%%%%%%%%%%%%%%%%%%%%%%%%%%%%%%%
%%% Slide-7
%%%%%%%%%%%%%%%%%%%%%%%%%%%%%%%%%%%%%%%%%%%%%%%%%%%%%%%%%%%%%%%%%%%%%%%%%


 \subsectiontitle{Arranging Boxes nicer}
       \subsection{\subsectitle}
       \begin{frame}
         \frametitle{\subsectitle}

        \begin{block}{Block 1}      The first block should be more up 
        \end{block}
\rule{0pt}{2cm}
        \begin{block}{Block 2}
          The second block should be more down
        \end{block}
      \end{frame}

%% %%%%%%%%%%%%%%%%%%%%%%%%%%%%%%%%%%%%%%%%%%%%%%%%%%%%%%%%%%%%%%%%%%%%%%%%%%%%

 \end{document}


